% GENERATED FILE. Edit canonical lessons, then run npm run generate:books.
% canonical-source-hash: fnv1a64:40ee9be2b0cdf821
% canonical-lessons: JA-C481-udemae, JA-C481-bukiyou, JA-C481-gijutsu, JA-C481-chie, JA-C481-kufuu-suru

\chapter{A Skilled Hand}
\label{ch:ja-a-skilled-hand}
\hlchaptermodality{481}

\begin{chapteropening}
\textbf{By the end of this chapter:} \emph{I can talk about skill, know-how and good ideas.}

Complete the last lesson of chapter 481: I can talk about skill, know-how and good ideas.
\end{chapteropening}

\section[\texorpdfstring{\ja{うでまえ} (udemae)}{うでまえ (udemae)}]{\ja{うでまえ} (udemae) — skill; ability}

\label{lesson:JA-C481-udemae}



\begin{quote}
Before the new one: say the Japanese for time and effort, then the Japanese for to finish off.
\end{quote}

\subsection*{You'll want to know: \ja{うでまえ}}
\textbf{\ja{うでまえ}} — \emph{udemae} — \textquotedblleft{}skill; ability\textquotedblright{}.

\begin{quote}
\textbf{\ja{あの} \ja{だいくは} \ja{うでまえが} \ja{いいです。}} — \emph{ano daiku wa udemae ga ii desu} — \textquotedblleft{}That carpenter is very skilled.\textquotedblright{}
\end{quote}

\begin{cousinweb}
\textbf{\ja{うでまえ}} (\emph{udemae}) is a native Japanese compound (wago), inherited rather than borrowed from Chinese: built from \emph{ude}, \textquotedblleft{}arm\textquotedblright{} and \emph{mae}, \textquotedblleft{}front\textquotedblright{} (JMdict, KANJIDIC2).
\end{cousinweb}

\begin{grammarlens}[title={What you've built}]
One more word for describing what you lived and felt.
\end{grammarlens}

\subsection*{Your turn}
\begin{itemize}
\raggedright
  \item \emph{Say it:} \emph{udemae}
  \item \emph{Say it:} \emph{udemae}, once more
  \item \emph{Say it:} \emph{ano daiku wa udemae ga ii desu} — That carpenter is very skilled.
  \item \emph{Recall:} say the Japanese for time and effort, then the Japanese for to finish off, then say \emph{udemae} again
\end{itemize}

\begin{tcolorbox}[breakable,colback=teal!4,colframe=teal!35!black,title={Before you move on}]
What does \ja{うでまえ} mean? (\textquotedblleft{}Skill; ability\textquotedblright{}.) Say it once more.
\end{tcolorbox}
\section[\texorpdfstring{\ja{ぶきよう} (bukiyō)}{ぶきよう (bukiyō)}]{\ja{ぶきよう} (bukiyō) — clumsy; awkward (with one's hands)}

\label{lesson:JA-C481-bukiyou}



\begin{quote}
Before the new one: say the Japanese for to finish off, then the Japanese for skill.
\end{quote}

\subsection*{You'll want to know: \ja{ぶきよう}}
\textbf{\ja{ぶきよう}} — \emph{bukiyō} — \textquotedblleft{}clumsy; awkward (with one's hands)\textquotedblright{}.

\begin{quote}
\textbf{\ja{あには} \ja{てが} \ja{ぶきようです。}} — \emph{ani wa te ga bukiyō desu} — \textquotedblleft{}My elder brother is clumsy with his hands.\textquotedblright{}
\end{quote}

A na-adjective: \ja{な} before a noun.

\begin{cousinweb}
\textbf{\ja{ぶきよう}} (\emph{bukiyō}) is Sino-Japanese (kango): three Chinese characters read with their on'yomi, readings borrowed from Chinese. They are \emph{bu}, \textquotedblleft{}negative\textquotedblright{}, \emph{ki}, \textquotedblleft{}utensil\textquotedblright{} (Middle Chinese \emph{khijH}) and \emph{yō}, \textquotedblleft{}utilize\textquotedblright{} (\emph{yowngH}) (KANJIDIC2; Middle Chinese in Baxter's transcription).
\end{cousinweb}

\begin{grammarlens}[title={What you've built}]
One more word for describing what you lived and felt.
\end{grammarlens}

\subsection*{Your turn}
\begin{itemize}
\raggedright
  \item \emph{Say it:} \emph{bukiyō}
  \item \emph{Say it:} \emph{bukiyō}, once more
  \item \emph{Say it:} \emph{ani wa te ga bukiyō desu} — My elder brother is clumsy with his hands.
  \item \emph{Recall:} say the Japanese for to finish off, then the Japanese for skill, then say \emph{bukiyō} again
\end{itemize}

\begin{tcolorbox}[breakable,colback=teal!4,colframe=teal!35!black,title={Before you move on}]
What does \ja{ぶきよう} mean? (\textquotedblleft{}Clumsy; awkward (with one's hands)\textquotedblright{}.) Say it once more.
\end{tcolorbox}
\section[\texorpdfstring{\ja{ぎじゅつ} (gijutsu)}{ぎじゅつ (gijutsu)}]{\ja{ぎじゅつ} (gijutsu) — technology; technique}

\label{lesson:JA-C481-gijutsu}



\begin{quote}
Before the new one: say the Japanese for skill, then the Japanese for clumsy.
\end{quote}

\subsection*{You'll want to know: \ja{ぎじゅつ}}
\textbf{\ja{ぎじゅつ}} — \emph{gijutsu} — \textquotedblleft{}technology; technique\textquotedblright{}.

\begin{quote}
\textbf{\ja{にほんの} \ja{ぎじゅつは} \ja{すごいです。}} — \emph{nihon no gijutsu wa sugoi desu} — \textquotedblleft{}Japanese technology is amazing.\textquotedblright{}
\end{quote}

\begin{cousinweb}
\textbf{\ja{ぎじゅつ}} (\emph{gijutsu}) is Sino-Japanese (kango): two Chinese characters read with their on'yomi, readings borrowed from Chinese. They are \emph{gi}, \textquotedblleft{}skill\textquotedblright{} (Middle Chinese \emph{gjeX}) and \emph{jutsu}, \textquotedblleft{}art\textquotedblright{} (\emph{zywit}) (KANJIDIC2; Middle Chinese in Baxter's transcription).
\end{cousinweb}

\begin{grammarlens}[title={What you've built}]
One more word for describing what you lived and felt.
\end{grammarlens}

\subsection*{Your turn}
\begin{itemize}
\raggedright
  \item \emph{Say it:} \emph{gijutsu}
  \item \emph{Say it:} \emph{gijutsu}, once more
  \item \emph{Say it:} \emph{nihon no gijutsu wa sugoi desu} — Japanese technology is amazing.
  \item \emph{Recall:} say the Japanese for skill, then the Japanese for clumsy, then say \emph{gijutsu} again
\end{itemize}

\begin{tcolorbox}[breakable,colback=teal!4,colframe=teal!35!black,title={Before you move on}]
What does \ja{ぎじゅつ} mean? (\textquotedblleft{}Technology; technique\textquotedblright{}.) Say it once more.
\end{tcolorbox}
\section[\texorpdfstring{\ja{ちえ} (chie)}{ちえ (chie)}]{\ja{ちえ} (chie) — wisdom; know-how}

\label{lesson:JA-C481-chie}



\begin{quote}
Before the new one: say the Japanese for clumsy, then the Japanese for technology.
\end{quote}

\subsection*{You'll want to know: \ja{ちえ}}
\textbf{\ja{ちえ}} — \emph{chie} — \textquotedblleft{}wisdom; know-how\textquotedblright{}.

\begin{quote}
\textbf{\ja{そふは} \ja{ちえが} \ja{あります。}} — \emph{sofu wa chie ga arimasu} — \textquotedblleft{}My grandfather is wise.\textquotedblright{}
\end{quote}

\begin{cousinweb}
\textbf{\ja{ちえ}} (\emph{chie}) is Sino-Japanese (kango): two Chinese characters read with their on'yomi, readings borrowed from Chinese. They are \emph{chi}, \textquotedblleft{}know\textquotedblright{} (Middle Chinese \emph{trje}) and \emph{e}, \textquotedblleft{}blessing\textquotedblright{} (\emph{hwejH}) (KANJIDIC2; Middle Chinese in Baxter's transcription).
\end{cousinweb}

\begin{grammarlens}[title={What you've built}]
One more word for describing what you lived and felt.
\end{grammarlens}

\subsection*{Your turn}
\begin{itemize}
\raggedright
  \item \emph{Say it:} \emph{chie}
  \item \emph{Say it:} \emph{chie}, once more
  \item \emph{Say it:} \emph{sofu wa chie ga arimasu} — My grandfather is wise.
  \item \emph{Recall:} say the Japanese for clumsy, then the Japanese for technology, then say \emph{chie} again
\end{itemize}

\begin{tcolorbox}[breakable,colback=teal!4,colframe=teal!35!black,title={Before you move on}]
What does \ja{ちえ} mean? (\textquotedblleft{}Wisdom; know-how\textquotedblright{}.) Say it once more.
\end{tcolorbox}
\section[\texorpdfstring{\ja{くふうする} (kufū suru)}{くふうする (kufū suru)}]{\ja{くふうする} (kufū suru) — to devise; to find a better way}

\label{lesson:JA-C481-kufuu-suru}



\begin{quote}
Before the new one: say the Japanese for technology, then the Japanese for wisdom.
\end{quote}

\subsection*{You'll want to know: \ja{くふうする}}
\textbf{\ja{くふうする}} — \emph{kufū suru} — \textquotedblleft{}to devise; to find a better way\textquotedblright{}.

\begin{quote}
\textbf{\ja{ははは} \ja{りょうりを} \ja{くふうします。}} — \emph{haha wa ryōri o kufūshimasu} — \textquotedblleft{}My mother keeps finding better ways to cook.\textquotedblright{}
\end{quote}

Polite \textbf{\ja{くふうします}}, \emph{kufūshimasu}; a \ja{する} verb. Past (polite): \textbf{\ja{ははが} \ja{りょうりを} \ja{くふうしました。}}, \emph{haha ga ryōri o kufūshimashita}.

\begin{cousinweb}
\textbf{\ja{くふうする}} (\emph{kufū suru}) is the Sino-Japanese (kango) noun \emph{kufū} with the native verb \emph{suru}, \textquotedblleft{}to do\textquotedblright{}: two Chinese characters read with their on'yomi, readings borrowed from Chinese. They are \emph{ku}, \textquotedblleft{}craft\textquotedblright{} (Middle Chinese \emph{kuwng}) and \emph{fū}, \textquotedblleft{}man\textquotedblright{} (\emph{pju}) (KANJIDIC2; Middle Chinese in Baxter's transcription).
\end{cousinweb}

\begin{grammarlens}[title={What you've built}]
One more word for describing what you lived and felt.
\end{grammarlens}

\subsection*{Your turn}
\begin{itemize}
\raggedright
  \item \emph{Say it:} \emph{kufū suru}
  \item \emph{Say it:} \emph{kufū suru}, once more
  \item \emph{Say it:} \emph{haha wa ryōri o kufūshimasu} — My mother keeps finding better ways to cook.
  \item \emph{Recall:} say the Japanese for technology, then the Japanese for wisdom, then say \emph{kufū suru} again
\end{itemize}

\begin{tcolorbox}[breakable,colback=teal!4,colframe=teal!35!black,title={Before you move on}]
What does \ja{くふうする} mean? (\textquotedblleft{}To devise; to find a better way\textquotedblright{}.) Say it once more.
\end{tcolorbox}
